%% toc: P1 8 8.1,8.2,8.5
\chapter{Readings of the Part 9 Program Constructors and Diamonds, the Autoepistemic Operator, the Part 19 Denial and Part 8 Operators in Their Presenting Chapters}
\chaptermark{Part 8, 9 and 19 Readings}
\label{sec:p01-w02c-p01-017-readings}
\index{precedence!program constructors}\index{precedence!autoepistemic operator}\index{propositional dynamic logic!bracketing}

\emph{Purpose.}  Several operators outside Part~8 have no rank,
or a rank without a fixity and association entry, in the precedence
apparatus of this part.  The program constructors \(\alpha;\beta\),
\(\alpha\cup\beta\), \(\alpha^{*}\) and \(A?\) of propositional dynamic
logic in Part~9 have the precedence entry \emph{n/a}; the grammar of
\cref{def:p09-pdl} generates programs without parentheses and so leaves
strings such as \(\alpha;\beta;\gamma\) and \(\alpha;\beta\cup\gamma\)
without a reading.  The diamond \(\langle\alpha\rangle X\), which the
grammar of Kozen and Parikh in \cref{sec:p09-presentation-n-c-021} takes
as primitive, needs a rank of its own, because that section prints it
without parentheses next to binary connectives, as in
\(\langle\alpha\rangle\bot\leftrightarrow\bot\) and
\(\langle\alpha^{*}\rangle X\to X\lor\langle\alpha^{*}\rangle(\neg
X\land\langle\alpha\rangle X)\).  The indexed diamond
\(\langle a\rangle p\) and box \([a]p\) of Kozen's modal
\(\mu\)-calculus in the model-checking section of Part~9
(\cref{sec:p09-n-c-018-kozen}) are printed in the same way, as in
\(\langle a\rangle p\land[a]q\).  The autoepistemic operator \(L\) of the
Gelfond--Lifschitz translation in Part~9 has formula rank~3: the prefix
reading specified here governs the unparenthesized occurrences in
\(\neg LB_1\land\cdots\land\neg LB_n\) in
\cref{def:p09-gl88-expansion}.  The denial \(\neg\) of the
selection-task section of Part~19 records the rank~60, which
Table~\ref{tab:p01-precedence} lists, but no table of this part gives its
fixity and association.  This chapter supplies a rank and an association,
or a reading fixed by delimiters or by the grammar, for each of them, and
it proves that the chosen association never changes the meaning of a
program.  Finally, the temporal, epistemic and linear operators of
Part~8 have ranks whose recorded scopes begin after the chapters that
define them; \cref{sec:p01-w02c-p01-017-p8-chapters} carries those ranks
into the defining chapters and into the annex and inquisitive passages
that print relational modalities, and it records where the grammar of the
classical linear presentation fixes a reading of its own.

\emph{Inputs.}  The fixity and association conventions
\cref{def:p01-w02-p01-026-precedence} and
\cref{def:p01-w02-p01-024-fixity}, with Table~\ref{tab:p01-precedence} and
Table~\ref{tab:p01-w02-p01-024-ranks}; the Part~8 program scale of item~20
of \cref{def:p01-w02b-p01-012-p8-readings}; the registry rows of the
operators, which record the program constructors, the box
\([\alpha]A\) and the diamond \(\langle\alpha\rangle A\) with the scope
\emph{P9 legacy unit 6 program-logics}, the
operator \(L\) with the scope \emph{P9 Gelfond-Lifschitz 1988
autoepistemic translation}, and the denial with the scope \emph{P19
sec:p19-r112-selection}; and the Part~9 and Part~19 passages that use
them: \cref{def:p09-pdl}, \cref{sec:p09-presentation-n-c-021},
\cref{def:p09-gl88-expansion} and \cref{sec:p19-r112-selection}; for
the last section, the chapters \cref{ch:modal-temporal-epistemic} and
\cref{ch:relevance-linear}, the annex proofs of
\cref{sec:p08-src-025-501}, the inquisitive chapter
\cref{sec:p08-inquisitive-propositional}, and the definition of linear
negation in the presentation of classical linear logic by
\textcite{Girard1995}.

\emph{Conclusions.}  \Cref{def:p01-w02c-p01-017-readings} places the
program constructors of Part~9 on a program scale (iteration first, then
composition, then choice, both binary constructors associating to the
left), gives the test a postfix reading with a delimiter rule for its
operand, ranks the diamond of propositional dynamic logic and the
indexed diamond and box of the modal \(\mu\)-calculus of the Part~9
model-checking section at~3 as prefix operators, ranks \(L\) at~3 as a
prefix operator, and gives the Part~19 denial the fixity \emph{prefix} with no association, at its recorded rank.
Table~\ref{tab:p01-w02c-p01-017-ranks} collects the entries.
\Cref{thm:p01-w02c-p01-017-pdl-bracketing} shows that every bracketing of
an iterated composition or choice of Part~9 programs has the same
transition relation, \cref{thm:p01-w02c-p01-017-prefix-chains} shows when
the rank of a prefix operator matters for a reading and when it does not,
and \cref{thm:p01-w02c-p01-017-printed-readings} derives the readings of
the formulas that the Part~9 and Part~19 passages print.
\Cref{def:p01-w02c-p01-017-p8-chapters} and
Table~\ref{tab:p01-w02c-p01-017-p8-chapter-ranks} give the Part~8
operators their ranks in their presenting chapters;
\cref{thm:p01-w02c-p01-017-chapter-readings} reads the formulas printed
there, and \cref{thm:p01-w02c-p01-017-cll-negation} shows that only the
source reading of the exponential De Morgan clauses defines linear
negation.  The proofs are in the annex.

\section{Assumptions and the grammar boundary}
\label{sec:p01-w02c-p01-017-assumptions}

The conventions of this chapter apply only after the language of the
passage has been fixed, as \cref{def:p01-w02-p01-024-fixity} requires,
and only inside the recorded scope of each row.  Where the grammar of a
passage fixes a reading, that reading stands.  The grammar of
\cref{def:p09-pdl} is
\[
 \alpha::=a\mid\alpha;\beta\mid\alpha\cup\beta\mid\alpha^{*}\mid A?,
 \qquad
 A::=p\mid\neg A\mid A\land B\mid[\alpha]A,
\]
a simultaneous definition of programs and formulas.  Read as a definition
of trees it is unambiguous, but as a definition of strings it is not: the
string \(a;b\cup c\) is generated both as a composition whose right
operand is \(b\cup c\) and as a choice whose left operand is \(a;b\).  The
same holds for the test-free grammar of Kozen and Parikh printed in
\cref{sec:p09-presentation-n-c-021}.  Neither passage states a bracketing
rule, so the convention below supplies one.  The box \([\alpha]A\)
delimits its program by the brackets, and its registry row records the
rank~3, so a program never stands next to a formula operator: inside the
brackets only program constructors and tests occur, and outside them the
box is one rank-3 prefix operator.  The diamond \(\langle\alpha\rangle X\)
delimits its program in the same way by the angle brackets.  In the
grammar of Kozen and Parikh it is primitive and the box is the defined
dual \([\alpha]X:=\neg\langle\alpha\rangle\neg X\); in the presentations
of Pratt and of Fischer and Ladner in the same section it is the dual
\(\neg[\alpha]\neg B\) of the box.  The formula operators of the scope are
the box, the diamond, the global negation (rank~3), conjunction (rank~4),
disjunction (rank~5), and the global implication and biconditional
(ranks~6 and~7).  The same scope also uses angle brackets for the
configuration \(\langle C,s\rangle\) of the while language, which pairs a
command with a state and stands only to the left of the evaluation sign
\(\Downarrow\); a diamond encloses one program and is followed by a
formula, so the two never share a position.

The language of \cref{def:p09-gl88-expansion} consists of Boolean formulas
and the operator \(L\), and the consequence operation \(th\) treats every
formula \(LF\) as a propositional atom of the Boolean skeleton.  The
translated rule
\((A_1\land\cdots\land A_m\land\neg LB_1\land\cdots\land\neg LB_n)\supset H\)
encloses its antecedent in parentheses, so the implication sign needs no
rank there; the reading of the antecedent depends on how \(L\), the
negation and the conjunction combine.  The negation in the translated
rule is the Boolean negation of the autoepistemic language, read by the
global negation row (rank~3); the negative body condition \(\neg B\) of
the program syntax is another row with the precedence \emph{n/a}, which
occurs only as a member of a comma-separated rule body.

In \cref{sec:p19-r112-selection} the letters \(p\) and \(q\) name
conditions, and the denial is written only directly before one of them,
as \(\neg q\) and \(\neg p\); combinations with a second condition are
stated in words (``\(p\) together with \(\neg q\)'').

\section{The convention}
\label{sec:p01-w02c-p01-017-convention}

\begin{hbconvention}{def:p01-w02c-p01-017-readings}{Readings of the Part~9 program constructors and diamonds, of the autoepistemic operator and of the Part~19 denial}
Each item holds only inside the recorded scope named with it.  Ranks of
formula operators are those of Table~\ref{tab:p01-precedence}, and a
smaller rank binds more tightly; the clauses of
\cref{def:p01-w02-p01-024-fixity} apply unchanged.
\begin{enumerate}
\item \emph{Test} \(A?\), scope \emph{P9 legacy unit 6 program-logics}:
  postfix; it forms a program from a formula.  Its operand is the longest
  formula that ends immediately before the sign \(?\).  Thus in \([a;p?]q\)
  the operand is \(p\), since \(a;p\) is not a formula, and in
  \([p\land q?]r\) it is \(p\land q\).
\item \emph{Iteration} \(\alpha^{*}\), same scope: postfix, program
  rank~1.  It applies to the immediately preceding atomic program,
  parenthesized program or test, so \(a;b^{*}\) is \(a;(b^{*})\) and
  \(p?^{*}\) is \((p?)^{*}\).
\item \emph{Composition} \(\alpha;\beta\), same scope: infix, program
  rank~2, left associative: \(\alpha;\beta;\gamma\) is
  \((\alpha;\beta);\gamma\).
\item \emph{Choice} \(\alpha\cup\beta\), same scope: infix, program
  rank~3, left associative: \(\alpha\cup\beta\cup\gamma\) is
  \((\alpha\cup\beta)\cup\gamma\).  Composition binds more tightly than
  choice, so \(\alpha;\beta\cup\gamma\) is \((\alpha;\beta)\cup\gamma\) and
  \(\alpha\cup\beta;\gamma\) is \(\alpha\cup(\beta;\gamma)\).
\item \emph{The program scale.}  The program ranks of items~2--4 are a
  scale of their own and are never compared with formula ranks.  A
  program occurs only between the brackets of a box, whose registry row
  records the formula rank~3, or as the operand of another program
  constructor; a formula occurs inside a program only as the operand of a
  test, which item~1 delimits.  The scale is the one that item~20 of
  \cref{def:p01-w02b-p01-012-p8-readings} sets up for the programs of
  \cref{def:p08-bvb501-pdl}, so the two dynamic logics of Parts~8 and~9
  read every program string alike.
\item \emph{Autoepistemic operator} \(L\), scope \emph{P9
  Gelfond-Lifschitz 1988 autoepistemic translation}: prefix, rank~3, no
  association (it is unary).  By clause~4 of
  \cref{def:p01-w02-p01-024-fixity} it takes the shortest following
  formula, and a chain of prefix signs nests: \(\neg LB\) is
  \(\neg(LB)\) and \(L\neg B\) is \(L(\neg B)\).  Against a binary
  connective it takes its operand first, so \(\neg LB_1\land\neg LB_2\) is
  \((\neg(LB_1))\land(\neg(LB_2))\).  This agrees with the source, which
  obtains the translation of a program by placing \(L\) directly after
  each negation sign, so that the negation applies to the modal atom
  \cite[p.~8]{GelfondLifschitz1988}.  When a formula metavariable stands
  as the operand, as in \(LF\) or \(\neg LF\), the operand is the whole
  formula that the metavariable denotes; written out for a compound
  formula, the instance carries parentheses, so for \(F=p\land q\) the
  formula \(LF\) is \(L(p\land q)\), not \((Lp)\land q\).
\item \emph{Denial} \(\neg\), scope \emph{P19 sec:p19-r112-selection}:
  prefix, at its recorded rank~60, no association (it is unary).  Its
  operand is the shortest following formula that clause~1 of
  \cref{def:p01-w02-p01-024-fixity} leaves to it, and a chain of denials
  nests.  Because the rank~60 exceeds the rank of every binary connective
  in Table~\ref{tab:p01-precedence}, a binary connective that follows the
  first letter of the operand takes that letter first; parentheses are
  therefore written whenever a denial in this scope is followed by a
  formula containing a binary connective, and a denied condition letter
  is written \(\neg p\) or \(\neg q\).
\item \emph{Diamond} \(\langle\alpha\rangle X\), scope \emph{P9 legacy
  unit 6 program-logics}: prefix, rank~3, no association (it is unary);
  its program is delimited by the angle brackets, inside which items~1--5
  apply.  By clause~1 of \cref{def:p01-w02-p01-024-fixity} it takes its
  operand before conjunction (rank~4), disjunction (rank~5), implication
  (rank~6) and the biconditional (rank~7), and by clause~4 a chain of
  diamonds, boxes and negations nests:
  \(\langle\alpha\rangle\langle\beta\rangle X\) is
  \(\langle\alpha\rangle(\langle\beta\rangle X)\) and
  \(\neg\langle\alpha\rangle\neg X\) is
  \(\neg(\langle\alpha\rangle(\neg X))\).  The rank is that of the box of
  the same scope, so the two dual modalities read alike, and it is the
  rank that the registry records for the indexed diamonds of Part~8
  (\cref{def:p08-bvb501-language}, and the program modalities of Kozen in
  \cref{sec:p08-conditionals-deontic-free}).
\item \emph{Diamond and box of the modal \(\mu\)-calculus}
  \(\langle a\rangle p\) and \([a]p\), indexed by a program constant
  \(a\), scope \emph{P9 legacy unit 5 model-checking}: prefix, rank~3,
  no association.  The diamond is primitive in the finitary language of
  Kozen in \cref{sec:p09-n-c-018-kozen} and the box is its defined dual
  \([a]p:=\neg\langle a\rangle\neg p\); \cref{def:p09-mu-calculus} closes
  its classes under both.  The rank is the one that the registry records
  for the same operators of the same source in Part~8 and for the
  fixed-point binders \(\mu X.\) and \(\nu X.\), whose scope includes
  Part~9.
\end{enumerate}
\end{hbconvention}

\begingroup\footnotesize
\begin{longtable}{@{}p{0.10\textwidth}p{0.17\textwidth}p{0.17\textwidth}p{0.20\textwidth}p{0.24\textwidth}@{}}
\caption{Ranks, fixity and association of the Part~9 program constructors and modalities, the autoepistemic operator and the Part~19 denial.}
\label{tab:p01-w02c-p01-017-ranks}\\
\toprule
Rank & Operator & Fixity & Association & Recorded scope \\
\midrule
\endfirsthead
\caption[]{Ranks, fixity and association of the Part~9 program constructors and modalities, the autoepistemic operator and the Part~19 denial (continued).}\\
\toprule
Rank & Operator & Fixity & Association & Recorded scope \\
\midrule
\endhead
\midrule\multicolumn{5}{r}{\emph{continued}}\\\endfoot
\bottomrule\endlastfoot
prog.~1 & \(\alpha^{*}\) & postfix & --- & P9 legacy unit 6 program-logics \\
prog.~2 & \(\alpha;\beta\) & infix & left & P9 legacy unit 6 program-logics \\
prog.~3 & \(\alpha\cup\beta\) & infix & left & P9 legacy unit 6 program-logics \\
--- & \(A?\) & postfix, formula to program & operand is the longest formula before \(?\) & P9 legacy unit 6 program-logics \\
3 & \([\alpha]A\) & prefix, program delimited by the brackets & --- & P9 legacy unit 6 program-logics \\
3 & \(\langle\alpha\rangle A\) & prefix, program delimited by the angle brackets & --- & P9 legacy unit 6 program-logics \\
3 & \(\langle a\rangle p\), \([a]p\) & prefix & --- & P9 legacy unit 5 model-checking \\
3 & \(L\) & prefix & --- & P9 Gelfond-Lifschitz 1988 autoepistemic translation \\
60 & \(\neg\) (denial of a condition) & prefix & --- & P19 sec:p19-r112-selection \\
\end{longtable}
\endgroup

\section{Results}
\label{sec:p01-w02c-p01-017-results}

The first result shows that the left association of items~3 and~4 selects
one program from a class of programs with the same meaning, so that it is
a choice of notation and not of semantics.  The semantic facts it uses are
those of \cref{def:p09-pdl}: composition, union and star denote relational
composition, union and reflexive-transitive closure, and composition
executes its left program first, as the identity
\(R_{\alpha^{*}}=I\cup(R_\alpha\mathbin{;}R_{\alpha^{*}})\) and the
computation \(R_{a;p?}=\{(s,t)\}\) of
\cref{sec:p09-presentation-n-c-021} use it.

\begin{hblemma}{thm:p01-w02c-p01-017-pdl-bracketing}{proved=annex:p01-w02c-p01-017-pdl-bracketing}{Bracketing does not affect the relation of an iterated Part~9 composition or choice}
Let \(M\) be a model of \cref{def:p09-pdl}, so that every program
\(\pi\) denotes a binary relation \(R_\pi\) on the states of \(M\), with
\(R_{\pi;\rho}=\{(w,v):(w,u)\in R_\pi\text{ and }(u,v)\in R_\rho\text{ for
some state }u\}\) and \(R_{\pi\cup\rho}=R_\pi\cup R_\rho\).  Let
\(n\geq1\) and let \(\pi_1,\ldots,\pi_n\) be programs.  For every
bracketing \(\gamma\) of \(\pi_1;\cdots;\pi_n\) and every bracketing
\(\kappa\) of \(\pi_1\cup\cdots\cup\pi_n\), each with the operands in this
order,
\begin{enumerate}
\item \((w,v)\in R_\gamma\) if and only if there are states
  \(u_0,\ldots,u_n\) with \(u_0=w\), \(u_n=v\) and
  \((u_{i-1},u_i)\in R_{\pi_i}\) for every \(i\) with \(1\leq i\leq n\);
\item \(R_\kappa=R_{\pi_1}\cup\cdots\cup R_{\pi_n}\).
\end{enumerate}
Hence any two bracketings of either iterated program have the same
relation, and \(M,s\models[\gamma]A\) and \(M,s\models[\kappa]A\) do not
depend on the bracketing, for every state \(s\) and formula \(A\).
\end{hblemma}

The second result separates the two situations in which a rank of a
prefix operator is or is not needed.  Write \(o\) for a prefix operator,
\(\star\) for a binary operator, and \(X,Y\) for operands that are atomic
or delimited, such as a letter, a parenthesized formula or a formula
metavariable.

\begin{hblemma}{thm:p01-w02c-p01-017-prefix-chains}{proved=annex:p01-w02c-p01-017-prefix-chains}{When the rank of a prefix operator matters}
Under the clauses of \cref{def:p01-w02-p01-024-fixity}:
\begin{enumerate}
\item A string \(o_1o_2\cdots o_kX\) with \(k\geq1\) prefix operators
  has the single reading \(o_1(o_2(\cdots(o_kX)\cdots))\), whatever the
  ranks of \(o_1,\ldots,o_k\), and also when some of them have no rank.
\item Let \(o\) have rank \(r\) and \(\star\) rank \(s\) with \(r\neq s\).
  The string \(oX\star Y\) has the single reading \((oX)\star Y\) when
  \(r<s\), and \(o(X\star Y)\) when \(r>s\).  When \(o\) has no rank, the
  string has no licensed reading.
\end{enumerate}
\end{hblemma}

The third result applies the convention to the formulas that the passages
print.  In the Part~9 items the global ranks of
\cref{def:p01-w02-p01-026-precedence} apply to \(\neg\), \(\land\),
\(\to\) and \(\leftrightarrow\).

\begin{hblemma}{thm:p01-w02c-p01-017-printed-readings}{proved=annex:p01-w02c-p01-017-printed-readings}{Readings of the printed Part~9 and Part~19 formulas}
Under \cref{def:p01-w02c-p01-017-readings}:
\begin{enumerate}
\item The box formulas \([a;p?]q\), \([p?;c]q\) and \([a;b^{*}]q\) of
  \cref{sec:p09-presentation-n-c-021} are \([a;(p?)]q\), \([(p?);c]q\) and
  \([a;(b^{*})]q\).
\item The star-unfolding, star-induction and test schemata of the PDL
  calculus,
  \([\alpha^{*}]A\leftrightarrow(A\land[\alpha][\alpha^{*}]A)\),
  \([\alpha^{*}](A\to[\alpha]A)\to(A\to[\alpha^{*}]A)\) and
  \([B?]A\leftrightarrow(B\to A)\), are
  \(([\alpha^{*}]A)\leftrightarrow(A\land([\alpha]([\alpha^{*}]A)))\),
  \(([\alpha^{*}](A\to([\alpha]A)))\to(A\to([\alpha^{*}]A))\) and
  \(([B?]A)\leftrightarrow(B\to A)\); and the induction principle
  \(X\land[\alpha^{*}](X\to[\alpha]X)\to[\alpha^{*}]X\) of
  \cref{sec:p09-presentation-n-c-021} is
  \((X\land([\alpha^{*}](X\to([\alpha]X))))\to([\alpha^{*}]X)\).
\item The translated antecedent
  \(A_1\land\cdots\land A_m\land\neg LB_1\land\cdots\land\neg LB_n\) of
  \cref{def:p09-gl88-expansion} is the left-associated conjunction of the
  \(m+n\) conjuncts \(A_1,\ldots,A_m,\neg(LB_1),\ldots,\neg(LB_n)\), and
  the members \(LF\) and \(\neg LF\) of \(\Delta(E)\) are \(L\) applied to
  \(F\) and the negation of that formula.
\item The denials \(\neg q\) and \(\neg p\) of
  \cref{sec:p19-r112-selection} have the same reading under the recorded
  rank~60 as under the rank~3 of the global negation.  For a string
  \(\neg X\land Y\) the two ranks give different readings: rank~60 gives
  \(\neg(X\land Y)\) and rank~3 gives \((\neg X)\land Y\).
\item In \cref{sec:p09-presentation-n-c-021} the definition
  \([\alpha]X:=\neg\langle\alpha\rangle\neg X\) defines \([\alpha]X\) as
  \(\neg(\langle\alpha\rangle(\neg X))\), and the axioms (2)--(7) of
  Kozen and Parikh are
  \begin{align*}
   &(\langle\alpha\rangle\bot)\leftrightarrow\bot,\\
   &(\langle\alpha\cup\beta\rangle X)\leftrightarrow
     ((\langle\alpha\rangle X)\lor(\langle\beta\rangle X)),\\
   &(\langle\alpha\rangle(X\lor Y))\leftrightarrow
     ((\langle\alpha\rangle X)\lor(\langle\alpha\rangle Y)),\\
   &(\langle\alpha;\beta\rangle X)\leftrightarrow
     (\langle\alpha\rangle(\langle\beta\rangle X)),\\
   &(\langle\alpha^{*}\rangle X)\leftrightarrow
     (X\lor(\langle\alpha\rangle(\langle\alpha^{*}\rangle X))),\\
   &(\langle\alpha^{*}\rangle X)\to
     (X\lor(\langle\alpha^{*}\rangle((\neg X)\land(\langle\alpha\rangle X)))).
  \end{align*}
  The formulas \(\langle a\rangle Q\lor\langle b\rangle Q\),
  \(\neg\langle a\rangle\neg Q\), \(C\land\langle a\rangle D\),
  \(Y\land\langle\beta\rangle\neg Y\) and
  \(\langle a\rangle A\to[a]A\) of the same section are
  \((\langle a\rangle Q)\lor(\langle b\rangle Q)\),
  \(\neg(\langle a\rangle(\neg Q))\), \(C\land(\langle a\rangle D)\),
  \(Y\land(\langle\beta\rangle(\neg Y))\) and
  \((\langle a\rangle A)\to([a]A)\).
\item In \cref{sec:p09-presentation-n-c-018} the formulas
  \(\langle a\rangle p\land[a]q\) and
  \(\mu Z.(\langle a\rangle Y\lor[b]Z)\) are
  \((\langle a\rangle p)\land([a]q)\) and
  \(\mu Z.((\langle a\rangle Y)\lor([b]Z))\), and the formula
  \(\neg\langle a\rangle\neg p\) that defines \([a]p\) is
  \(\neg(\langle a\rangle(\neg p))\).
\end{enumerate}
\end{hblemma}

\section{Worked example}
\label{sec:p01-w02c-p01-017-example}

\begin{hbexample}{def:p01-w02c-p01-017-example}{Reading and evaluating a Part~9 program}
Take the three-state model of \cref{sec:p09-presentation-n-c-021}:
states \(s,t,u\), \(R_a=\{(s,t),(s,u)\}\) and \(V(p)=V(q)=\{t\}\), and add
an atomic program \(b\) with \(R_b=\{(t,u)\}\).
\begin{enumerate}
\item By items~1, 3 and~4 the string \(a;p?;b\cup a\) is
  \(((a;(p?));b)\cup a\).  The test relation is \(R_{p?}=\{(t,t)\}\), so
  \(R_{a;p?}=\{(s,t)\}\) and \(R_{(a;p?);b}=\{(s,u)\}\); the choice adds
  \(R_a\), giving \(\{(s,t),(s,u)\}\).  Hence
  \(s\models[a;p?;b\cup a]q\) fails, because \(u\) is a successor and
  \(u\notin V(q)\).
\item The other bracketing \(a;((p?;b)\cup a)\) of the same string, which
  the convention does not license, has the relation \(\{(s,u)\}\): the
  program \((p?;b)\cup a\) has the relation
  \(\{(t,u)\}\cup R_a=\{(t,u),(s,t),(s,u)\}\); from \(s\) the program \(a\)
  reaches \(t\) and \(u\), from \(t\) the second program reaches \(u\),
  and from \(u\) it reaches nothing.  The two relations differ, so the
  relative rank of composition and choice decides the meaning of the
  string.
\item By \cref{thm:p01-w02c-p01-017-pdl-bracketing} the strings
  \(a;p?;b\) and \(a;(p?;b)\) have the same relation \(\{(s,u)\}\), so the
  left association of item~3 changes nothing there.
\item In the Gelfond--Lifschitz translation of the rule
  \(p\leftarrow q,\neg r,\neg s\), the antecedent
  \(q\land\neg Lr\land\neg Ls\) is \((q\land\neg(Lr))\land\neg(Ls)\), a
  conjunction of an atom and two negated modal atoms, each of which
  \(th\) treats as a negated propositional letter.
\item In the model of item~1, the string \(\langle a\rangle p\lor q\)
  is \((\langle a\rangle p)\lor q\) by item~8.  At the state \(t\) it is
  true, because \(t\in V(q)\).  The reading \(\langle a\rangle(p\lor q)\),
  which a diamond ranked above disjunction would give, is false at \(t\),
  because \(t\) has no \(a\)-successor.  So the rank of the diamond
  against disjunction decides the truth value of the string at \(t\).
\end{enumerate}
\end{hbexample}

\section{Failure cases and limitations}
\label{sec:p01-w02c-p01-017-limits}

\begin{hbwarning}{def:p01-w02c-p01-017-limits}{Limits of the readings}
\begin{enumerate}
\item The ranks of items~2--4 are a convention of this part, not a rule
  stated in \cref{def:p09-pdl} or in \cref{sec:p09-presentation-n-c-021};
  the printed Part~9 formulas write parentheses wherever composition meets
  choice, as in \([(a;b^{*})\cup(p?;c)]q\), so none of them depends on
  it.  As the worked example shows, the relative rank of composition and
  choice changes the relation of a program, whereas the association of
  either one alone does not (\cref{thm:p01-w02c-p01-017-pdl-bracketing}).
\item The rank~3 of the diamonds in items~8 and~9 is a convention of
  this part; neither \cref{sec:p09-presentation-n-c-021} nor
  \cref{sec:p09-presentation-n-c-018} states a bracketing rule.  Every
  printed Part~9 formula listed in
  \cref{thm:p01-w02c-p01-017-printed-readings} keeps its reading under
  any rank below~4, since a diamond there meets other prefix operators,
  where ranks do not matter (\cref{thm:p01-w02c-p01-017-prefix-chains}),
  or binary connectives of rank at least~4.  A rank above that of a
  binary connective would change readings: the soundness argument for
  axiom (7) in \cref{sec:p09-presentation-n-c-021} needs the antecedent
  \(\langle\alpha^{*}\rangle X\) and the consequent
  \(X\lor\langle\alpha^{*}\rangle(\neg X\land\langle\alpha\rangle X)\),
  which only a diamond binding more tightly than \(\to\) and \(\lor\)
  delivers, and item~5 of the worked example gives a separating case.
\item The sign \(*\) has two roles in the same scope: the postfix
  iteration \(\alpha^{*}\), set as a superscript on a program, and the
  separating conjunction \(P*Q\) of rank~4, set on the line between two
  assertions.  The position and the sort of the operands tell them apart;
  neither is ever an operand of the other.
\item The rank~3 of item~6 holds only for \(L\) in its recorded scope.
  The letters \(L_i\) of the rule bodies in the same section are rule
  literals, not operators, and the belief operator \(LA\) of the Part~8
  further-reading section is another row with its own scope, although it
  also records the rank~3.
\item The recorded rank~60 of the Part~19 denial lies above every binary
  rank of Table~\ref{tab:p01-precedence}, so it reverses the reading that
  the global negation gives to an unparenthesized \(\neg X\land Y\)
  (\cref{thm:p01-w02c-p01-017-printed-readings}).  Item~7 avoids the
  conflict by writing parentheses; the section itself prints no such
  string.  The value is recorded in the registry, and a rank~3 for this
  row would give the same reading to every formula the section prints.
\end{enumerate}
\end{hbwarning}

\section{Part 8 operators in the chapters that present them}
\label{sec:p01-w02c-p01-017-p8-chapters}
\index{precedence!temporal operators}\index{precedence!epistemic operators}\index{precedence!linear connectives}\index{linear negation!reading of the De Morgan clauses}

The ranks of the temporal next, until, eventually and always operators,
of the knowledge, mutual-knowledge, common-knowledge, belief and
distributed-knowledge operators, and of the branching-time,
fixed-point, tense, announcement and satisfaction operators in
Table~\ref{tab:p01-w02-p01-024-ranks} carry recorded scopes that begin
with the chapter of conditionals, deontic logic and free logic
(\cref{sec:p08-conditionals-deontic-free}).  The chapter that defines
these operators, \cref{ch:modal-temporal-epistemic}, comes before it: it
gives the language of linear temporal logic
(\cref{sec:modal-temporal-epistemic-02}), its Hilbert calculus with the
until-unfolding schema, the completeness proof of
\cref{sec:p08-linear-temporal-completeness}, branching time
(\cref{def:p08-ctl-semantics,def:p08-ctl-star-semantics}), Kozen's
\(\mu\)-calculus (\cref{def:p08-modal-mu-semantics}), Kamp's tense
connectives (\cref{def:p08-tense-until-since}), multi-agent epistemic
logic (\cref{sec:modal-temporal-epistemic-03}), public announcement
(\cref{def:p08-public-announcement}) and basic hybrid logic
(\cref{def:p08-hybrid-k}).  In the same way the ranks of the linear
connectives carry the scope of the chapter on paraconsistent and
nonmonotonic logics (\cref{sec:p08-paraconsistent-and-nonmonotonic}),
which uses none of them, whereas they are defined in
\cref{ch:relevance-linear}: multiplicative intuitionistic linear logic
in \cref{sec:relevance-linear-03} and \cref{entry:p08-imll}, classical
linear logic in \cref{sec:p08-cll} and the Lambek calculus in
\cref{sec:p08-lambek}.  The relational presentation of
\cref{sec:p08-src-025-501} ranks its modalities and its satisfaction
operator for the body section (items~3 and~10 of
\cref{def:p01-w02b-p01-012-p8-readings}), but its annex proofs print the
same operators outside that section, and the inquisitive chapter prints
the relational box once, in \(\Box p\lor\Box\neg p\)
(\cref{sec:p08-inquisitive-propositional}).

\Cref{def:p01-w02c-p01-017-p8-chapters} gives every one of these
operators, in the scope of the passage that presents or prints it, the
rank, fixity and association it has in
Table~\ref{tab:p01-w02-p01-024-ranks}, and it records where the grammar
of a presentation fixes a reading on its own.  The scope names are the
reader regions of the registry: \emph{P8 modal-temporal-epistemic} is
\cref{ch:modal-temporal-epistemic} with the annex items of its
completeness and decision proofs, and \emph{P8 relevance-linear} is
\cref{ch:relevance-linear} with the annex item of the multiplicative
calculus.

\begin{hbconvention}{def:p01-w02c-p01-017-p8-chapters}{Ranks of Part~8 temporal, epistemic, modal and linear operators in their presenting chapters}
Each item holds inside the scope named with it, after the grammar of the
passage has been fixed.  Ranks are those of
Table~\ref{tab:p01-precedence}; the clauses of
\cref{def:p01-w02-p01-024-fixity} apply unchanged, so two different
operators of the same rank never compete and are separated by
parentheses.  Negation has rank~3, conjunction~4, disjunction~5,
implication~6 (right associative) and the biconditional~7
(nonassociative).
\begin{enumerate}
\item \emph{Linear temporal operators}, scope \emph{P8
  modal-temporal-epistemic}: next \(\mathsf X\), eventually \(\mathsf F\)
  and always \(\mathsf G\) are prefix operators of rank~3; until
  \(A\mathbin{\mathsf U}B\) is infix of rank~3.5 and nonassociative.  The
  grammar of \cref{sec:modal-temporal-epistemic-02} writes every binary
  construction in parentheses; the ranks govern the schemas, rules,
  definitions and examples that omit them.
\item \emph{Branching-time operators}, same scope: the path quantifiers
  \(\mathsf E\) and \(\mathsf A\) and the paired operators
  \(\mathsf{EX}\), \(\mathsf{AX}\), \(\mathsf{EF}\), \(\mathsf{AF}\) are
  prefix operators of rank~3.  In the state formulas
  \(\mathsf E(A\mathbin{\mathsf U}B)\) and
  \(\mathsf A(A\mathbin{\mathsf U}B)\) of
  \cref{def:p08-ctl-semantics} the parentheses are part of the grammar:
  a path quantifier stands before a path formula, and an until formula
  is one only with its parentheses.
\item \emph{Fixed points and Kozen's modalities}, same scope: the
  binders \(\mu Z.\) and \(\nu Z.\) of \cref{def:p08-modal-mu-semantics}
  are rank-3 prefix binders, which by clause~5 of
  \cref{def:p01-w02-p01-024-fixity} take the shortest formula after the
  dot, since that definition states no rule for the extent of the body;
  the diamond \(\langle a\rangle A\) and the defined box \([a]A\) are
  prefix operators of rank~3.
\item \emph{Tense operators}, same scope: Kamp's \(U\,A\,B\) and
  \(S\,A\,B\) of \cref{def:p08-tense-until-since} are prefix operators of
  rank~3 with two operands, each operand being the shortest formula in
  its position, and Prior's past operators \(\mathsf P\) and
  \(\mathsf H\) are prefix operators of rank~3.  Kamp's prefix \(U\) and
  the infix until of item~1 are different operators.
\item \emph{Epistemic operators}, same scope: knowledge \(K_i\),
  mutual knowledge \(E_G\), common knowledge \(C_G\), belief \(B_i\) and
  distributed knowledge \(D_G\) are prefix operators of rank~3.  In the
  announcement modalities \([!\chi]A\) and \(\langle!\chi\rangle A\) the
  brackets delimit the announced formula \(\chi\) completely, and the box
  or diamond so formed is a prefix operator of rank~3.  The satisfaction
  operator \(@_i\) of \cref{def:p08-hybrid-k} is a prefix operator of
  rank~3 whose operand is the shortest formula after the nominal
  subscript.
\item \emph{Relational modalities in the annex proofs of
  \cref{sec:p08-src-025-501}}, scope \emph{the annex items from
  \cref{annex:p08.src025501.iso} through
  \cref{annex:p08.src025501.topology-def}}: the indexed diamond and box
  \(\langle m\rangle\varphi\) and \([m]\varphi\), the single-modality
  abbreviations \(\Diamond\) and \(\Box\) and the satisfaction operator
  \(@_i\) are prefix operators of rank~3, as items~3 and~10 of
  \cref{def:p01-w02b-p01-012-p8-readings} give them in the body section.
\item \emph{Relational box}: in the scope \emph{P8
  inquisitive-propositional}, which is the inquisitive chapter with its
  annex, the box \(\Box\) is a prefix operator of rank~3.
\item \emph{Linear connectives}, scope \emph{P8 relevance-linear}: the
  exponentials \(!\) and \(?\) are prefix of rank~2.5 and linear
  negation \({}^\perp\) is postfix of rank~2.5; tensor \(\otimes\) and
  par \(\bindnasrepma\) are infix of rank~4, left associative; plus
  \(\oplus\) is infix of rank~5, left associative; linear implication
  \(\multimap\) is infix of rank~6, right associative.  The Lambek
  product \(A\mathbin{\cdot}B\) is infix of rank~4, left associative, and
  the divisions \(A\backslash B\) and \(B/A\) are infix of rank~6 and
  nonassociative.  The grammar of \cref{sec:relevance-linear-03},
  \(A::=p\mid\mathbf 1\mid A\otimes A\mid A\multimap A\), and the language
  of \cref{def:p08-lambek-semantics} generate strings without
  parentheses and state no bracketing rule, so these ranks supply one.
\item \emph{With}, scope \cref{sec:p08-cll}: the additive conjunction
  \(A\mathbin{\&}B\) is infix of rank~4, left associative.  The strong
  conjunction of the t-norm logics, written with the same sign and also
  of rank~4, has a recorded scope that includes this chapter, but it is
  not used in \cref{sec:p08-cll}: in the definition of the language, the
  connective table and the phase interpretation of that section the sign
  is always the additive conjunction.
\item \emph{Grammar boundary for linear negation}, scope
  \cref{def:p08-cll-language}: in the displayed De Morgan clauses that
  define linear negation, the right-hand sides \(?A^\perp\) and
  \(!A^\perp\) are read \(?(A^\perp)\) and \(!(A^\perp)\), and every
  \({}^\perp\) on a right-hand side applies to the immediately preceding
  letter or parenthesized formula.  This is the reading of the
  presentation that the definition follows, in which linear negation is
  \emph{defined} by these equations, clause by clause on the outermost
  connective \cite[\S1.2.2, p.~10]{Girard1995}.  Inside those clauses the
  grammar of the definition replaces clause~3 of
  \cref{def:p01-w02-p01-024-fixity}; everywhere else in the scope a
  prefix exponential next to the postfix negation is written with
  parentheses, \(!(A^\perp)\) or \((!A)^\perp\).
\end{enumerate}
\end{hbconvention}

\begingroup\footnotesize
\begin{longtable}{@{}p{0.06\textwidth}p{0.24\textwidth}p{0.12\textwidth}p{0.135\textwidth}p{0.27\textwidth}@{}}
\caption{Ranks of Part~8 operators in the chapters that present or print them.}
\label{tab:p01-w02c-p01-017-p8-chapter-ranks}\\
\toprule
Rank & Operators & Fixity & Association & Scope \\
\midrule
\endfirsthead
\caption[]{Ranks of Part~8 operators in the chapters that present or print them (continued).}\\
\toprule
Rank & Operators & Fixity & Association & Scope \\
\midrule
\endhead
\midrule\multicolumn{5}{r}{\emph{continued}}\\\endfoot
\bottomrule\endlastfoot
2.5 & \(!A\), \(?A\) & prefix & --- & P8 relevance-linear \\
2.5 & \(A^\perp\) & postfix & --- & P8 relevance-linear; source reading in \cref{def:p08-cll-language} \\
3 & \(\mathsf X,\mathsf F,\mathsf G\) & prefix & --- & P8 modal-temporal-epistemic \\
3 & \(\mathsf E,\mathsf A;\ \mathsf{EX},\mathsf{AX},\mathsf{EF},\mathsf{AF}\) & prefix & --- & P8 modal-temporal-epistemic \\
3 & \(\mu Z.A,\nu Z.A\) & prefix binder & --- & P8 modal-temporal-epistemic \\
3 & \(\langle a\rangle A,[a]A\) & prefix & --- & P8 modal-temporal-epistemic \\
3 & \(U\,A\,B,\ S\,A\,B\) (Kamp) & prefix, two operands & --- & P8 modal-temporal-epistemic \\
3 & \(\mathsf P,\mathsf H\) & prefix & --- & P8 modal-temporal-epistemic \\
3 & \(K_i,E_G,C_G,B_i,D_G\) & prefix & --- & P8 modal-temporal-epistemic \\
3 & \([!\chi]A,\langle!\chi\rangle A\); \(@_iA\) & prefix & --- & P8 modal-temporal-epistemic \\
3 & \(\langle m\rangle\varphi,[m]\varphi\); \(\Box,\Diamond\); \(@_i\varphi\) & prefix & --- & annex items of \cref{sec:p08-src-025-501} \\
3 & \(\Box\) & prefix & --- & P8 inquisitive-propositional \\
3.5 & \(A\mathbin{\mathsf U}B\) & infix & none & P8 modal-temporal-epistemic \\
4 & \(\otimes\), \(\bindnasrepma\) & infix & left & P8 relevance-linear \\
4 & \(A\mathbin{\&}B\) (with) & infix & left & \cref{sec:p08-cll} \\
4 & \(A\mathbin{\cdot}B\) (Lambek product) & infix & left & P8 relevance-linear \\
5 & \(A\oplus B\) & infix & left & P8 relevance-linear \\
6 & \(\multimap\) & infix & right & P8 relevance-linear \\
6 & \(A\backslash B,\ B/A\) (Lambek divisions) & infix & none & P8 relevance-linear \\
\end{longtable}
\endgroup

The first result applies the convention to the formulas that these
passages print.  The second shows that the reading of item~10 is the only
one under which the De Morgan clauses define an operation on formulas.

\begin{hblemma}{thm:p01-w02c-p01-017-chapter-readings}{proved=annex:p01-w02c-p01-017-chapter-readings}{Readings of printed formulas of the Part~8 presenting chapters}
Under \cref{def:p01-w02c-p01-017-p8-chapters}:
\begin{enumerate}
\item The schemas and the induction rule of the Hilbert calculus for
  linear temporal logic in \cref{ch:modal-temporal-epistemic} read
  \begin{align*}
   &(\mathsf X(A\to B))\to((\mathsf XA)\to(\mathsf XB)),\qquad
    (\mathsf X(\neg A))\leftrightarrow(\neg(\mathsf XA)),\\
   &(A\mathbin{\mathsf U}B)\leftrightarrow
     (B\lor(A\land(\mathsf X(A\mathbin{\mathsf U}B)))),\\
   &\text{premise }(B\lor(A\land(\mathsf XC)))\to C,\qquad
    \text{conclusion }(A\mathbin{\mathsf U}B)\to C,
  \end{align*}
  and the definitions \(\mathsf FA\defeq\topv\mathbin{\mathsf U}A\) and
  \(\mathsf GA\defeq\neg\mathsf F\neg A\) define \(\mathsf FA\) as
  \(\topv\mathbin{\mathsf U}A\) and \(\mathsf GA\) as
  \(\neg(\mathsf F(\neg A))\).
\item The schemas of the multi-agent calculus read
  \begin{align*}
   &(K_i(A\to B))\to((K_iA)\to(K_iB)),\qquad(K_iA)\to A,\\
   &(K_iA)\to(K_i(K_iA)),\qquad(\neg(K_iA))\to(K_i(\neg(K_iA)));
  \end{align*}
  the common-knowledge fixed-point axiom \(C_GA\to E_G(A\land C_GA)\)
  reads \((C_GA)\to(E_G(A\land(C_GA)))\); and the announced formula
  \(p\land\neg K_ip\) of \cref{def:p08-unsuccessful-announcement} reads
  \(p\land(\neg(K_ip))\).
\item The branching-time definition
  \(\mathsf E\alpha\defeq\neg\mathsf A\neg\alpha\) defines
  \(\mathsf E\alpha\) as \(\neg(\mathsf A(\neg\alpha))\), and the
  temporal formula \(\mathsf G\mathsf Fp\) of the worked diagrams is
  \(\mathsf G(\mathsf Fp)\).
\item In the annex proof \cref{annex:p08.src025501.e-collapse}, the formula
  \(\Diamond\Diamond p\to\Box(\Diamond p\lor\Box p)\) reads
  \((\Diamond(\Diamond p))\to(\Box((\Diamond p)\lor(\Box p)))\), and in the
  inquisitive chapter \(\Box p\lor\Box\neg p\) reads
  \((\Box p)\lor(\Box(\neg p))\).
\item The tensor clause of \cref{def:p08-cll-language} has the
  right-hand side \(A^\perp\bindnasrepma B^\perp\), which reads
  \((A^\perp)\bindnasrepma(B^\perp)\), and the definiens of
  \(A\multimap B\defeq A^\perp\bindnasrepma B\) reads
  \((A^\perp)\bindnasrepma B\).  In the scope \emph{P8
  relevance-\allowbreak{}linear} the string \(A\multimap B\multimap C\) is
  \(A\multimap(B\multimap C)\), and \(A\otimes B\multimap C\) is
  \((A\otimes B)\multimap C\).
\end{enumerate}
\end{hblemma}

\begin{hblemma}{thm:p01-w02c-p01-017-cll-negation}{proved=annex:p01-w02c-p01-017-cll-negation}{Only the source reading of the exponential clauses defines linear negation}
Let the CLL formulas be the trees generated from the paired atoms
\(p,p^\perp\) and the units \(\mathbf 1,\bot,\top,0\) by
\(\otimes,\bindnasrepma,\mathbin{\&},\oplus,!,?\), as in
\cref{def:p08-cll-language}, so that every formula is an atom, a unit,
or has exactly one outermost constructor and immediate subformulas.
\begin{enumerate}
\item Read the clauses as in item~10 of
  \cref{def:p01-w02c-p01-017-p8-chapters}: \(p\mapsto p^\perp\),
  \(p^\perp\mapsto p\), \(\mathbf 1\mapsto\bot\), \(\bot\mapsto\mathbf 1\),
  \(\top\mapsto0\), \(0\mapsto\top\), and for the compound formulas
  \((A\otimes B)^\perp=A^\perp\bindnasrepma B^\perp\),
  \((A\bindnasrepma B)^\perp=A^\perp\otimes B^\perp\),
  \((A\mathbin{\&}B)^\perp=A^\perp\oplus B^\perp\),
  \((A\oplus B)^\perp=A^\perp\mathbin{\&}B^\perp\),
  \((!A)^\perp=?(A^\perp)\) and \((?A)^\perp=!(A^\perp)\).  Then there is
  exactly one operation on formulas satisfying these equations, and it
  is an involution: \((A^\perp)^\perp=A\) for every formula \(A\).
\item Read the two exponential clauses instead as the equations
  \((!A)^\perp=(?A)^\perp\) and \((?A)^\perp=(!A)^\perp\).  Then no
  involution on formulas satisfies them, and the equations together do
  not determine the value of \((!A)^\perp\) for any formula \(A\).
\end{enumerate}
\end{hblemma}

\begin{hbexample}{def:p01-w02c-p01-017-p8-example}{Ranks that decide a value}
\begin{enumerate}
\item Take the temporal structure with \(V(1)=\{p\}\), \(V(2)=\{q\}\) and
  \(V(n)=\varnothing\) for every other \(n\).  By item~1 of
  \cref{def:p01-w02c-p01-017-p8-chapters} the string
  \(\mathsf Xp\mathbin{\mathsf U}q\) is
  \((\mathsf Xp)\mathbin{\mathsf U}q\), since next (3) takes \(p\) before
  until (3.5).  At position~0 it is false: \(q\) holds only at
  position~2, and \(\mathsf Xp\) fails at position~1 because \(p\) fails
  at position~2.  The other bracketing
  \(\mathsf X(p\mathbin{\mathsf U}q)\) is true at position~0, because at
  position~1 the formula \(p\) holds and \(q\) holds at position~2.  So
  the rank of next against until decides the truth value of the string.
\item Take an agent \(j\), worlds \(w,v\), the relation \(\sim_j\) the
  full relation on \(\{w,v\}\), and \(p\) true only at \(w\).  By item~5
  the string \(K_jp\lor p\) is \((K_jp)\lor p\), which is true at \(w\)
  because \(p\) is.  The reading \(K_j(p\lor p)\), which a knowledge
  operator ranked above disjunction would give, is false at \(w\),
  because \(v\) is a \(j\)-alternative of \(w\) at which \(p\) fails.
\item In \cref{def:p08-cll-language}, linear negation of
  \(A\otimes!B\) is computed by the tensor clause and then the bang
  clause, each read as in item~10:
  \((A\otimes!B)^\perp=A^\perp\bindnasrepma(!B)^\perp
  =A^\perp\bindnasrepma?(B^\perp)\).  Under the other reading of the bang
  clause the last step would give \((?B)^\perp\), the negation of a
  why-not formula, in place of the why-not formula \(?(B^\perp)\);
  \cref{thm:p01-w02c-p01-017-cll-negation} shows that this reading does
  not define linear negation at all.
\end{enumerate}
\end{hbexample}

\begin{hbwarning}{def:p01-w02c-p01-017-p8-limits}{What the chapter readings do not decide}
\begin{enumerate}
\item The items give the operators the ranks that
  Table~\ref{tab:p01-w02-p01-024-ranks} records for them in later
  scopes; they introduce no new rank, and the presenting chapters state
  no bracketing rule of their own except where item~2 and item~10 say
  so.  Different operators of the same rank, such as \(\otimes\) and
  \(\bindnasrepma\), or \(!\) and \({}^\perp\) outside the De Morgan
  clauses, still need parentheses.
\item A chain of implications written without parentheses is one formula,
  read from the right.  The annex proof
  \cref{annex:p08.src025501.e-collapse} prints
  \(\Diamond p\to\Box\Diamond p\to\Diamond\Box p\to\Box p\) and
  \(\Box p\to p\to\Diamond p\) for the successive implications between
  adjacent members, each of which it derives separately; the convention
  reads the first display as
  \(\Diamond p\to(\Box\Diamond p\to(\Diamond\Box p\to\Box p))\), so the
  display is to be taken as a list of the three adjacent implications,
  not as an abbreviation licensed by the ranks.
\item The sign \(\&\) names two connectives whose recorded scopes
  include \cref{ch:relevance-linear} (item~9).  The reading is fixed by
  the section, not by the rank, which is~4 for both.
\item The boundary of item~10 covers only the displayed De Morgan clauses
  of \cref{def:p08-cll-language}.  It does not license \(!A^\perp\) or
  \(?A^\perp\) as abbreviations elsewhere, where
  \cref{def:p01-w02-p01-024-fixity} requires parentheses.
\end{enumerate}
\end{hbwarning}

\section{Relationships}
\label{sec:p01-w02c-p01-017-relationships}

The convention \cref{def:p01-w02c-p01-017-readings} extends
\cref{def:p01-w02-p01-024-fixity} and
\cref{def:p01-w02b-p01-012-p8-readings} to registry rows of Parts~9
and~19: the program scale of items~2--5 is the scale of item~20 of the
Part~8 convention, applied to the scope of \cref{def:p09-pdl};
\cref{thm:p01-w02c-p01-017-pdl-bracketing} is the Part~9 counterpart of
\cref{thm:p01-w02b-p01-012-program-bracketing}; and the ranks of items~6
to~9 sit beside the prefix operators of rank~3 in
Table~\ref{tab:p01-w02-p01-024-ranks}.  The diamonds of items~8 and~9
receive the rank that Table~\ref{tab:p01-precedence} lists for the box
of the same Part~9 scope, for the indexed diamonds of Part~8, and for
the fixed-point binders whose scope runs through Part~9.  The box \([\alpha]A\), the
separating conjunction \(P*Q\) and the other ranked rows of the scope
\emph{P9 legacy unit 6 program-logics} keep the entries of that table.
The registry extracts of \cref{sec:p01-w02b-p01-012-extracts} list the
rows treated here with their recorded scopes and precedence entries.

\Cref{def:p01-w02c-p01-017-p8-chapters} adds no rank: it repeats, for
the scopes of the chapters and passages that present or print the
Part~8 operators, the ranks, fixities and associations that
Table~\ref{tab:p01-w02-p01-024-ranks} and
\cref{def:p01-w02-p01-026-precedence} give the same operators in the
later scopes, and its item~6 extends items~3 and~10 of
\cref{def:p01-w02b-p01-012-p8-readings} from the body section of the
relational presentation to its annex proofs.  The worked reading of the
until-unfolding schema in \cref{def:p01-w02-p01-024-rank-examples} and
the linear-logic examples there therefore hold in the chapters that
state those schemas and connectives.  Item~10 marks the one place where a
presentation's own grammar replaces clause~3 of
\cref{def:p01-w02-p01-024-fixity}, which the example
\(!A^\perp\) of that clause otherwise excludes.

\emph{Sources.}  The readings rest on the handbook passages cited with
each item: the grammar and relational semantics of propositional dynamic
logic (\cref{def:p09-pdl}), the presentations of Kozen and Parikh and of
Harel, Kozen and Tiuryn with the three-state model
(\cref{sec:p09-presentation-n-c-021}), Kozen's modal \(\mu\)-calculus
in the model-checking section (\cref{sec:p09-n-c-018-kozen},
\cref{def:p09-mu-calculus}), the stable expansion and program
translation of Gelfond and Lifschitz (\cref{def:p09-gl88-expansion}; the
translation itself on page~8,
\cite[p.~8]{GelfondLifschitz1988}), and
the selection-task section (\cref{sec:p19-r112-selection}), together with
the registry rows \texttt{registry/symbols.csv} of the program
constructors, the box, the diamonds, the operator \(L\) and the denial.
The Part~8 readings rest on the defining passages named in
\cref{sec:p01-w02c-p01-017-p8-chapters}: the language, calculus and
completeness proof of linear temporal logic, the branching-time, fixed-point,
tense, epistemic, announcement and hybrid definitions of
\cref{ch:modal-temporal-epistemic}, the annex proofs of
\cref{sec:p08-src-025-501}, the inquisitive chapter, and the linear
calculi of \cref{ch:relevance-linear}.  The reading of the exponential
De Morgan clauses is that of Girard's presentation, where linear
negation is defined by those equations
\cite[\S1.2.2, p.~10]{Girard1995}.
